% =====================================================================
% Cover letter — Magnetic Resonance Imaging (Elsevier)
% Built with tectonic: `tectonic cover_letter.tex`
% =====================================================================
\documentclass[11pt]{letter}
\usepackage[margin=1in]{geometry}
\usepackage{microtype}
\usepackage[hidelinks]{hyperref}

\signature{Avery Karlin, BA \\ Department of Applied Physics and Applied Mathematics,
Columbia University \\ Department of Computer Science, University of Colorado Boulder \\
\texttt{ak5232@columbia.edu}}

\address{Avery Karlin \\ Columbia University \\ New York, NY, USA}

\begin{document}

\begin{letter}{Professor John C.\ Gore \\ Editor-in-Chief, \textit{Magnetic Resonance Imaging}
\\ Elsevier}

\opening{Dear Professor Gore,}

I am pleased to submit our manuscript, ``Calibration and Efficiency of Uncertainty Estimates
in Intravoxel Incoherent Motion Imaging: Quantile Intervals, Cross-Paradigm Comparison, and a
Cram\'er--Rao Audit of Amortized Posteriors,'' for consideration as an original research
article in \textit{Magnetic Resonance Imaging}.

The work asks a question that matters wherever quantitative MR parameters carry error bars:
when an IVIM fitting method reports an uncertainty, can it be believed? Using a single
biexponential IVIM simulation framework, the study makes two methodological contributions.
First, it compares uncertainty constructions across classical, Bayesian, and deep-learning
paradigms and identifies which yield calibrated intervals on the bounded, skewed parameters
$D^{*}$ and $f$ --- showing that quantile intervals, not Gaussian credible intervals, recover
nominal coverage. Second, and more generally, it benchmarks an amortized neural posterior
estimator against the Cram\'er--Rao information floor, demonstrating that a posterior can pass
aggregate calibration checks while being overconfident and biased exactly where the parameter
is weakly identified. The transferable message is a diagnostic one: information-floor (CRLB)
auditing, paired with shape-aware intervals, exposes pointwise miscalibration that aggregate
coverage conceals --- a check applicable to learned quantitative-MRI posteriors well beyond
IVIM.

We believe the manuscript is well matched to the computational and methodological scope of
\textit{Magnetic Resonance Imaging}. It is squarely a methods contribution to MR parameter
estimation and uncertainty quantification: it concerns simulation-based inference, the
information-theoretic limits of MR fitting, and deployment-time safeguards (an
out-of-distribution gate and per-scheme recalibration) for learned estimators. All code, the
calibration grid, and the efficiency map are released openly as a reproducibility module and
offered to the OSIPI initiative, consistent with the journal's emphasis on rigorous,
reproducible MR methodology.

This manuscript is original, has not been published previously, and is not under consideration
for publication elsewhere. Its submission is approved by the author, and there are no competing
interests to declare. The author takes full responsibility for the integrity of the work; a
declaration of the use of AI-assisted writing tools is included with the manuscript.

% [[DECISION: disclose prior MRM review? author's call]]
% This letter deliberately makes NO mention of any prior review history at another journal.
% Whether to proactively disclose the earlier Magnetic Resonance in Medicine review to the
% MRI editor is the author's decision; if desired, add a short, factual paragraph here.

Thank you for considering our work. I would be glad to suggest suitable reviewers on request
and to respond to any questions.

\closing{Sincerely,}

\end{letter}
\end{document}
